%%%%%%%%%%%%%%%%%%%%%%%%%%%%%%%%%%%%%%%%%%%%%%%%%%%%%%%%%%%%%%%%%%%%%%%%%%%%%%
%
% Dirac Delta
%
%%%%%%%%%%%%%%%%%%%%%%%%%%%%%%%%%%%%%%%%%%%%%%%%%%%%%%%%%%%%%%%%%%%%%%%%%%%%%%

\chapter{Dirac Delta}

%%%%%%%%%%%%%%%%%%%%%%%%%%%%%%%%%%%%%%%%%%%%%%%%%%%%%%%%%%%%%%%%%%%%%%%%%%%%%%
\section{Definition}

The Dirac delta is "heuristically" defined by

\[
\delta(x)=
\begin{cases}
0,&x\neq0\\
\infty,&x=0
\end{cases}
\]
together with the normalization

\[
\int_{-\infty}^{+\infty}\delta(x)\,dx=1.
\]
Its defining property is the \textit{sifting} property

\[
\boxed{
\int_{-\infty}^{+\infty}
f(x)\,\delta(x-x_0)\,dx
=
f(x_0)
}
\]
which, by a change of variable, can be also interpreted as a convolution integral:
\[
\boxed{
\int_{-\infty}^{+\infty}
f(\alpha)\,\delta(x-\alpha)\,d\alpha
= f(x) * \delta(x) = f(x).
}
\]
Thus, the Dirac delta is the identity element of the convolution.
%%%%%%%%%%%%%%%%%%%%%%%%%%%%%%%%%%%%%%%%%%%%%%%%%%%%%%%%%%%%%%%%%%%%%%%%%%%%%%
\section{Basic Properties}

\begin{center}
\begin{tabular}{ll}
\toprule
Property & Relation\\
\midrule

Sifting &
$\int f(x)\delta(x-x_0)\,dx=f(x_0)$
\\[0.3cm]

Even symmetry &
$\delta(-x)=\delta(x)$
\\[0.3cm]

Scaling &
$\displaystyle
\delta(ax)=\frac{1}{|a|}\delta(x)
$
\\[0.3cm]

Derivative &
$\displaystyle
\int f(x)\delta'(x-x_0)\,dx=-f'(x_0)
$
\\[0.3cm]

Product (sampling) &
$\displaystyle
f(x)\delta(x-x_0)
=
f(x_0)\delta(x-x_0)
$
\\%[0.3cm]

% Composition &
% $\displaystyle
% \delta(f(x))
% =
% \sum_i
% \frac{\delta(x-x_i)}
% {|f'(x_i)|}
% $
%\\
\bottomrule
\end{tabular}
\end{center}

% where \(x_i\) are the simple zeros of \(f(x)\).


%%%%%%%%%%%%%%%%%%%%%%%%%%%%%%%%%%%%%%%%%%%%%%%%%%%%%%%%%%%%%%%%%%%%%%%%%%%%%%
\section{Two-Dimensional Delta}

\[
\delta(\rvec)
=
\delta(x)\delta(y)
\]

with

\[
\iint
\delta(\rvec)\,d\rvec
=
1.
\]

More generally,

\[
\delta(\rvec-\rvec_0)
=
\delta(x-x_0)\delta(y-y_0).
\]
The sifting property becomes:
\[
\boxed{
\iint
f(x,y)
\delta(x-x_0)
\delta(y-y_0)
\,dx\,dy
=
f(x_0,y_0)
}
\]
More generally:
\[
\boxed{
\iint
f(\rvec)
\delta(\rvec-\rvec_0)
\,d\rvec
=
f(\rvec_0)
}
\]


%%%%%%%%%%%%%%%%%%%%%%%%%%%%%%%%%%%%%%%%%%%%%%%%%%%%%%%%%%%%%%%%%%%%%%%%%%%%%%
\section{Delta as a Limit}

The Dirac delta can be understood heuristically as the limit of a sequence
of increasingly narrow functions with unit area.

For example,

\[
\boxed{
\delta(x)
=
\lim_{a\rightarrow\infty}
a\,\sinc(ax)
}
\]

where

\[
\sinc(x)=\frac{\sin(\pi x)}{\pi x}.
\]

% Indeed,

% \[
% \int_{-\infty}^{+\infty}
% a\,\sinc(ax)\,dx
% =
% 1.
% \]

Similarly, a Gaussian representation is

\[
\boxed{
\delta(x)
=
\lim_{\sigma\rightarrow0}
\frac{1}{\sqrt{2\pi}\sigma}
e^{-x^2/(2\sigma^2)}
}
\]

%%%%%%%%%%%%%%%%%%%%%%%%%%%%%%%%%%%%%%%%%%%%%%%%%%%%%%%%%%%%%%%%%%%%%%%%%%%%%%
\section{Delta from Complex Exponentials}

The Dirac delta can also be represented as

\[
\boxed{
\delta(x)
=
\int_{-\infty}^{+\infty}
e^{i2\pi f x}\,df
}
\]

Equivalently,

\[
\boxed{
\delta(x)
=
\lim_{B\rightarrow\infty}
\int_{-B}^{B}
e^{i2\pi f x}\,df
}
\]

By solving the integral we obtain

\[
\delta(x)
=
\lim_{B\rightarrow\infty}
2B\,\sinc(2Bx),
\]
which is consistent with the limit discussed in the section above.
%%%%%%%%%%%%%%%%%%%%%%%%%%%%%%%%%%%%%%%%%%%%%%%%%%%%%%%%%%%%%%%%%%%%%%%%%%%%%%
\section{Fourier Transform}

The fundamental Fourier-transform pair is

\[
\boxed{
\delta(x)
\xleftrightarrow{\ \FT\ }
1
}
\]

and, in two dimensions,

\[
\boxed{
\delta(x)\delta(y)
\xleftrightarrow{\ \FT\ }
1
}
\]

Conversely,

\[
\boxed{
1
\xleftrightarrow{\ \FT\ }
\delta(f)
}
\]

in one dimension, and

\[
\boxed{
1
\xleftrightarrow{\ \FT\ }
\delta(f_x)\delta(f_y)
}
\]

in two dimensions.